\section{Objectifs}
\begin{itemize}
\item Simuler une variable aléatoire.
\item Calculer la moyenne d’un échantillon.
\item Comparer expérimentalement moyenne observée et espérance.
\end{itemize}

\section{Prérequis}
Variables aléatoires, moyenne, fonctions et boucles Python.

\section{Activité de découverte}
Simuler de nombreux échantillons permet d’observer la fluctuation de leur moyenne autour de l’espérance.

\section{Vocabulaire}
Simulation, échantillon, taille, moyenne d’échantillon, espérance.

\section{Cours}
Un échantillon de taille $n$ est une liste de $n$ réalisations d’une variable aléatoire. Sa moyenne vaut
\[
m=\frac{x_1+\cdots+x_n}{n}.
\]

Le BO demande d’étudier expérimentalement la distance entre $m$ et l’espérance $\mu$, puis de simuler $N$ échantillons et d’observer la proportion des cas où
\[
|m-\mu|\le \frac{2\sigma}{\sqrt n}.
\]

Une simulation illustre un modèle ; elle ne constitue pas une preuve.

\section{Exemples corrigés}
Pour un dé équilibré, une fonction qui renvoie un entier de 1 à 6 simule une réalisation. Une liste de $n$ appels forme un échantillon.

\section{Méthode}
Découper le programme en fonctions : simuler une valeur, construire un échantillon, calculer sa moyenne, répéter l’expérience.

\section{Erreurs fréquentes}
\begin{itemize}
\item Confondre fréquence observée et probabilité exacte.
\item Confondre $N$, nombre d’échantillons, et $n$, taille de chaque échantillon.
\item Modifier le modèle entre deux séries que l’on compare.
\end{itemize}

\section{Synthèse}
La simulation permet de relier modèle probabiliste, échantillonnage et comportement des moyennes observées.
